\documentclass[11pt]{article}
\usepackage[margin=1in]{geometry}
\usepackage{amsmath,amssymb}
\usepackage{lmodern}
\usepackage[T1]{fontenc}
\usepackage{enumitem}
\usepackage{microtype}

\title{2012 F=ma Exam Review}
\author{Nathan P.}
\date{December 26, 2025}

\begin{document}
\maketitle

\section*{Problem 1}
Trick: write $\Delta x$ for both the drop hitting the sink and the one in the air, and divide them.

\section*{Problem 2}
Use the kinematics formulas to get the angle.

\section*{Problem 8}
Don't forget to write little things like $mg$.

\section*{Problem 10}
The time depends on $a_{cm}$, which only depends on $\beta$ for $I=\beta mr^2$ when rolling without slipping.

\section*{Problem 14}
It's simpler if you just do a torque balance with the negative-mass hole; no need to find $x_{cm}$.
Also, careful: ``weight'' means the number already includes $g$.

\section*{Problem 17}
Log-log is a power graph. (Just take the log of both sides of the function to check how it's linearized. If the simplified form leaves $\log y$ and $\log x$, then it'd only be linear on a log-log graph.)
For $ax^b$, the slope is the power $b$ and the intercept is $\log a$.

\section*{Problem 18}
The period is independent of the acceleration of the system.
(This problem assumes the box is much more massive, so it isn't accelerated by the spring. Dropping a box with an oscillating mass inside makes it a two-body oscillator with a higher frequency, as in 2017 \#2. The period is independent of frame, so for the oscillation to not change, $M_{\text{box}}\gg m$.)
You can tell the amplitude, and therefore the max speed, changes because the equilibrium position shifts.
The height of max speed and the max height also change because the equilibrium shifts.

\section*{Problem 19}
\[
P=\frac{\Delta E}{t}=\frac{K+U}{t}=\frac mt\left(\tfrac12v^2+gh\right)=\rho Av\left(\tfrac12v^2+gh\right)=\rho\pi r^2v\left(\tfrac12v^2+gh\right).
\]
Solve for $r$, using $dm/dt=\rho Av$.

\section*{Problem 23}
You cannot find $g$ just from the inclined plane angle, masses, and time; you need the distance travelled. (Just think: $g$ has units of m/s$^2$, but there's nothing providing the meters. Also, $mgh$ would need $h$, which is a distance.)

\section*{Problem 24}
Pay attention to the stretch: it's $\ell$, not $2\ell$.
Use triangle geometry, like the $\tfrac23\cdot\tfrac{\sqrt3}{2}$ side.

\end{document}
